\section{Discussion}

The results support a chemometric interpretation of molecular benchmark scores: predictive performance is jointly determined by the model, the dataset, and the validation partition. Random, ordinary scaffold, and target-balanced scaffold protocols changed target distributions, chemical proximity, scaffold composition, and predictive gaps. A score without these accompanying properties does not fully specify the validation problem that was measured.

The null audit showed that the target-balanced search performed its intended statistical role. In five datasets, it produced substantially smaller target gaps than ordinary scaffold splitting and than most random scaffold assignments of comparable size. This matters because it demonstrates that the counterfactual was not merely a different random scaffold split. The search systematically controlled a measurable source of partition shift while preserving scaffold disjointness. HIV provided an important boundary case: its ordinary scaffold partition was already well balanced, so further first-moment optimization offered little statistical leverage.

The predictive findings also show why target balance should not be equated with benchmark validity. Reducing target-mean mismatch yielded smaller generalization gaps for nine dataset--model combinations, but larger gaps for six. ClinTox and ESOL contained within-dataset direction reversals across models. These reversals imply that a partition change can interact with model inductive bias, feature use, and the selected chemical region. The balanced partition does not hold every other property constant; it controls one objective while the chemical composition of the test set changes. It should therefore be interpreted as a diagnostic intervention rather than an improved leaderboard split.

FreeSolv illustrates a second limitation of partition design. Its target gap decreased strongly, yet the selected test size remained farther from the nominal 20\% target because a small number of large scaffold groups constrained feasible assignments. This is not a search failure that can be removed by stronger claims. It is evidence that some datasets cannot simultaneously provide exact test size, scaffold disjointness, target balance, and stable sample support. Such structural constraints should be reported rather than hidden behind a single aggregate score.

Random partitions remain useful when the intended use case resembles interpolation among chemically related compounds. Their generally higher maximum test-to-train Tanimoto similarity, however, means that they should not be used as unqualified evidence of transfer to new chemical series. Ordinary scaffold partitions answer a more distant chemical question, but their target shift and scaffold concentration can confound that question. The practical solution is not to designate one universal split. It is to define the intended validation population and report the properties of the resulting partition.

For molecular property prediction studies, a minimum reporting set should include the exact partition algorithm, random seeds, train and test target summaries, test-size deviation, scaffold overlap and concentration, test-to-train chemical similarity, uncertainty across repeated partitions or model fits, and model-ranking stability. When a split is optimized using target information, the objective and constraints should be stated explicitly and its behaviour should be compared with a target-blind reference distribution.

\section{Limitations}

This study has several limitations. First, the predictive models are fingerprint-based classical baselines rather than graph neural networks or molecular transformers. This choice reduces representation-related variability and makes repeated validation inexpensive, but future work should test whether deep molecular encoders show the same partition interactions. Second, the target-balanced search controls only test size and the first moment of the target distribution. It does not match variances, tails, class-conditional chemical space, or other covariates, and it is not guaranteed to find the global optimum. Third, the benchmark contains six datasets. The 20 seeds quantify robustness within these fixed datasets and should not be treated as 20 independent chemical populations. Fourth, the ordinary scaffold partition is deterministic, so its repeated results contain model-level but not partition-level variation. Fifth, the chemical-similarity audit uses fingerprint-based Tanimoto proximity and cannot represent every aspect of chemical or mechanistic novelty. Finally, no temporal or prospective validation was available.

\section{Conclusion}

Sample partitioning should be treated as a measured component of chemometric model validation. Target-balanced scaffold search can reduce response-distribution mismatch far beyond target-blind random scaffold assignment, but its effect on predictive performance is heterogeneous and model-dependent. Random partitions, ordinary scaffold partitions, and target-balanced scaffold partitions each combine different levels of chemical proximity, target shift, and scaffold concentration. Reporting these diagnostics alongside model scores provides a more transparent basis for judging molecular generalization than any single split or leaderboard value alone.